\chapter{Boolean Analysis — Low Degree}
%%% generated by the blueprint pipeline from TCSlib.BooleanAnalysis.BLR.LowDegree
\section{Overview}

This module develops the theory of low-degree testing on the Boolean hypercube
$\{0,1\}^n$.  Starting from the XOR group structure, it defines the Gowers product
(multiplicative derivative product), the Gowers uniformity norms, and the
Reed--Muller code $\mathrm{RM}(d,n)$ as the set of Boolean functions of degree
$\le d$.  The main results are a Fourier-analytic characterisation of degree
($\|\cdot\|_{U^{d+1}} = 1 \Leftrightarrow \deg \le d$), completeness and soundness
of the $(d{+}1)$-fold derivative test, and a quantitative soundness bound showing
that if $f$ is $\varepsilon$-far from $\mathrm{RM}(d,n)$ then the test rejects with
probability at least $\varepsilon$.

%%%%%%%%%%%%%%%%%%%%%%%%%%%%%%%%%%%%%%%%%%%%%%%%%%%%%%%%%%%%%%%%%%%%%%%%%%%%%%%
\section{Declarations}

\subsection{Hypercube algebra}

\begin{lemma}[Commutativity of componentwise XOR]
\lean{LowDegreeTest.xor_vec_comm}
\label{LowDegreeTest.xor_vec_comm}
\leanok
\uses{BoolFourier.hypercube, BoolFourier.xor_vec}
\difficulty{1}
For all $x, y \in \{0,1\}^n$, the componentwise XOR is symmetric: the vectors with
$i$-th coordinates $x_i \oplus y_i$ and $y_i \oplus x_i$ coincide, so $x \oplus y = y
\oplus x$.
\end{lemma}

\begin{lemma}[XOR with the all-zeros vector]
\lean{LowDegreeTest.xor_vec_zero}
\label{LowDegreeTest.xor_vec_zero}
\leanok
\uses{BoolFourier.hypercube, BoolFourier.xor_vec, BoolFourier.zero_vec}
\difficulty{1}
For every $x \in \{0,1\}^n$, the componentwise XOR of $x$ with the all-zeros vector
$\mathbf{0}$ equals $x$; that is, $x \oplus \mathbf{0} = x$. Equivalently, $\mathbf{0}$
is a right identity for the group operation on $(\bbf_2)^n$.
\end{lemma}

\begin{lemma}[XOR of a vector with itself]
\lean{LowDegreeTest.xor_vec_self}
\label{LowDegreeTest.xor_vec_self}
\leanok
\uses{BoolFourier.hypercube, BoolFourier.xor_vec, BoolFourier.zero_vec}
\difficulty{2}
For every $x \in \{0,1\}^n$, the componentwise XOR of $x$ with itself is the all-zeros
vector: $x \oplus x = \mathbf{0}$.
\end{lemma}

\begin{lemma}[Associativity of componentwise XOR]
\lean{LowDegreeTest.xor_vec_assoc}
\label{LowDegreeTest.xor_vec_assoc}
\leanok
\uses{BoolFourier.hypercube, BoolFourier.xor_vec}
\difficulty{1}
For all $x, y, z \in \{0,1\}^n$, the componentwise XOR is associative: writing $x \oplus
y$ for the vector with $i$th coordinate $x_i \oplus y_i$,
\[
(x \oplus y) \oplus z = x \oplus (y \oplus z).
\]
\end{lemma}

\subsection{Gowers product}

\begin{definition}[Multiplicative derivative]
\lean{LowDegreeTest.mult_deriv}
\label{LowDegreeTest.mult_deriv}
\leanok
\statementsource{boolean-ch06-pseudorandomness-f2}{pdf-6c10bfcc8080-p019-b005}
\uses{BoolFourier.BoolFun, BoolFourier.hypercube, BoolFourier.xor_vec}
For a function $f : \{0,1\}^n \to \bbr$ and a direction vector $h \in \{0,1\}^n$,
the multiplicative derivative is defined by
$(\Delta_h f)(x) = f(x) \cdot f(x \oplus h)$.
\end{definition}

\begin{definition}[Gowers product]
\lean{LowDegreeTest.gowers_product}
\label{LowDegreeTest.gowers_product}
\leanok
\statementsource{boolean-ch06-pseudorandomness-f2}{pdf-6c10bfcc8080-p026-b004}
\uses{BoolFourier.BoolFun, BoolFourier.hypercube, BoolFourier.xor_vec}
For $f : \{0,1\}^n \to \bbr$, base point $x \in \{0,1\}^n$, and direction
vectors $h_1, \ldots, h_k \in \{0,1\}^n$, the order-$k$ Gowers product is defined
recursively:
\[
  \mathrm{GP}(f, 0, x, \emptyset) = f(x),\quad
  \mathrm{GP}(f, k{+}1, x, h_1, \ldots, h_{k+1})
  = \mathrm{GP}(f, k, x, h_1,\ldots,h_k)
    \cdot \mathrm{GP}(f, k, x \oplus h_{k+1}, h_1,\ldots,h_k).
\]
It averages $f$ over the $2^k$ vertices of the combinatorial cube
$x + \mathrm{span}\{h_1,\ldots,h_k\}$.
\end{definition}

\begin{lemma}[Base case of the Gowers product]
\lean{LowDegreeTest.gowers_product_zero}
\label{LowDegreeTest.gowers_product_zero}
\leanok
\uses{BoolFourier.BoolFun, BoolFourier.hypercube, LowDegreeTest.gowers_product}
\difficulty{1}
Let $f : \{0,1\}^n \to \bbr$ be a Boolean function and let $x \in \{0,1\}^n$. With no
direction vectors supplied, the order-$0$ Gowers product equals the value of $f$ at $x$:
\[
  \mathrm{GP}(f, 0, x, \emptyset) = f(x).
\]
\end{lemma}

\begin{lemma}[Recursive unfolding of the Gowers product]
\lean{LowDegreeTest.gowers_product_succ}
\label{LowDegreeTest.gowers_product_succ}
\leanok
\uses{BoolFourier.BoolFun, BoolFourier.hypercube, BoolFourier.xor_vec, LowDegreeTest.gowers_product}
\difficulty{1}
Let $f : \{0,1\}^n \to \bbr$ be a Boolean function, let $k \ge 0$, let $x \in \{0,1\}^n$
be a base point, and let $h_1, \ldots, h_{k+1} \in \{0,1\}^n$ be direction vectors. Then
the order-$(k{+}1)$ Gowers product splits along its last direction as
\[
  \mathrm{GP}(f, k{+}1, x, h_1, \ldots, h_{k+1})
  = \mathrm{GP}(f, k, x, h_1, \ldots, h_k)
    \cdot \mathrm{GP}(f, k, x \oplus h_{k+1}, h_1, \ldots, h_k),
\]
where $x \oplus h_{k+1}$ denotes the componentwise XOR.
\end{lemma}

\begin{lemma}[Multiplicative derivative as the order-one Gowers product]
\lean{LowDegreeTest.mult_deriv_eq_gowers_one}
\label{LowDegreeTest.mult_deriv_eq_gowers_one}
\leanok
\uses{BoolFourier.BoolFun, BoolFourier.hypercube, LowDegreeTest.gowers_product, LowDegreeTest.mult_deriv}
\difficulty{1}
Let $f : \{0,1\}^n \to \bbr$ be a Boolean function, and let $x, h \in \{0,1\}^n$. Then
the multiplicative derivative of $f$ in direction $h$, evaluated at $x$, coincides with
the order-$1$ Gowers product of $f$ at base point $x$ with the single direction vector
$h$; that is,
\[
  (\Delta_h f)(x) \;=\; \mathrm{GP}(f, 1, x, h) \;=\; f(x)\, f(x \oplus h),
\]
where $\oplus$ denotes componentwise XOR.
\end{lemma}

\begin{lemma}[Gowers products of a $\pm1$-valued function take values $\pm1$]
\lean{LowDegreeTest.gowers_product_pm1}
\label{LowDegreeTest.gowers_product_pm1}
\leanok
\uses{BoolBLR.lift_pm1, BoolFourier.hypercube, BoolFourier.xor_vec, LowDegreeTest.gowers_product, LowDegreeTest.gowers_product_succ}
\difficulty{4}
Let $f : \{0,1\}^n \to \{0,1\}$ be a Boolean-valued function, and let $F : \{0,1\}^n \to
\bbr$, $F(x) = (-1)^{f(x)}$, be its $\pm1$ lift. Then for every order $k \ge 0$, every
base point $x \in \{0,1\}^n$, and all direction vectors $h_1, \ldots, h_k \in
\{0,1\}^n$, the order-$k$ Gowers product satisfies
\[
  \mathrm{GP}(F, k, x, h_1, \ldots, h_k) \in \{1, -1\}.
\]
\end{lemma}

\subsection{Multi-expectation}

\begin{definition}[Multi-expectation]
\lean{LowDegreeTest.multi_expectation}
\label{LowDegreeTest.multi_expectation}
\leanok
\uses{BoolFourier.hypercube}
For $g : (\{0,1\}^n)^k \to \bbr$, the multi-expectation is the uniform average
\[
  \E_{h_1,\ldots,h_k}[g(h_1,\ldots,h_k)]
  = \frac{1}{2^{nk}} \sum_{h_1,\ldots,h_k \in \{0,1\}^n} g(h_1,\ldots,h_k).
\]
\end{definition}

\begin{lemma}[Cardinality of $k$-tuples over the Boolean hypercube]
\lean{LowDegreeTest.card_multi_hypercube}
\label{LowDegreeTest.card_multi_hypercube}
\leanok
\uses{BoolFourier.card_hypercube, BoolFourier.hypercube}
\difficulty{2}
Fix natural numbers $n$ and $k$, and let $\{0,1\}^n$ denote the Boolean hypercube of
dimension $n$. The set of $k$-tuples of points of $\{0,1\}^n$, that is
$\left(\{0,1\}^n\right)^k$, has exactly $2^{nk}$ elements.
\end{lemma}

\begin{lemma}[Multi-expectation of a constant]
\lean{LowDegreeTest.multi_expectation_const}
\label{LowDegreeTest.multi_expectation_const}
\leanok
\uses{BoolFourier.card_hypercube, BoolFourier.hypercube, LowDegreeTest.multi_expectation}
\difficulty{3}
For any constant $c \in \bbr$ and any $k$, the uniform average over $k$ points
$h_1,\ldots,h_k$ of the Boolean hypercube $\{0,1\}^n$ of the constant function equal to
$c$ is itself $c$:
\[
  \E_{h_1,\ldots,h_k}[c] = c.
\]
\end{lemma}

\subsection{Polynomial degree}

\begin{definition}[Degree $\le d$ for $\pm 1$ functions]
\lean{LowDegreeTest.is_degree_le_pm1}
\label{LowDegreeTest.is_degree_le_pm1}
\leanok
\statementsource{boolean-ch06-pseudorandomness-f2}{pdf-6c10bfcc8080-p007-b004}
\uses{BoolFourier.BoolFun, BoolFourier.hypercube, LowDegreeTest.gowers_product}
A function $f : \{0,1\}^n \to \bbr$ has degree $\le d$ if all its $(d{+}1)$-fold
multiplicative derivatives vanish: for every $x$ and every $h_1,\ldots,h_{d+1}$,
$\mathrm{GP}(f, d{+}1, x, h_1,\ldots,h_{d+1}) = 1$.
\end{definition}

\begin{definition}[Degree $\le d$ for Boolean functions]
\lean{LowDegreeTest.is_degree_le_bool}
\label{LowDegreeTest.is_degree_le_bool}
\leanok
\statementsource{boolean-ch06-pseudorandomness-f2}{pdf-6c10bfcc8080-p007-b004}
\uses{BoolBLR.lift_pm1, BoolFourier.hypercube, LowDegreeTest.is_degree_le_pm1}
A Boolean function $f : \{0,1\}^n \to \{0,1\}$ has degree $\le d$ if its $\pm 1$
lift $(-1)^f$ has degree $\le d$ in the sense of \texttt{LowDegreeTest.is\_degree\_le\_pm1}.
\end{definition}

\begin{lemma}[Degree bound is monotone]
\lean{LowDegreeTest.degree_le_succ}
\label{LowDegreeTest.degree_le_succ}
\leanok
\uses{BoolFourier.BoolFun, BoolFourier.xor_vec, LowDegreeTest.is_degree_le_pm1}
\difficulty{4}
Let $f : \{0,1\}^n \to \bbr$, and recall that $f$ has degree $\le d$ when all of its
$(d{+}1)$-fold multiplicative derivatives vanish, i.e. the Gowers product
$\mathrm{GP}(f, d{+}1, x, h_1,\ldots,h_{d+1}) = 1$ for every base point $x \in
\{0,1\}^n$ and all directions $h_1,\ldots,h_{d+1} \in \{0,1\}^n$. If $f$ has degree $\le
d$, then $f$ also has degree $\le d+1$; that is, $\mathrm{GP}(f, d{+}2, x,
h_1,\ldots,h_{d+2}) = 1$ for every $x$ and all directions $h_1,\ldots,h_{d+2} \in
\{0,1\}^n$.
\end{lemma}

\begin{lemma}[Degree at most zero characterizes constant functions]
\lean{LowDegreeTest.degree_le_zero_iff_constant}
\label{LowDegreeTest.degree_le_zero_iff_constant}
\leanok
\uses{BoolBLR.lift_pm1, BoolFourier.BoolToPM1, BoolFourier.hypercube, BoolFourier.xor_vec, LowDegreeTest.is_degree_le_bool}
\difficulty{5}
A Boolean function $f : \{0,1\}^n \to \{0,1\}$ has degree $\le 0$ if and only if it is
constant, that is, $f(x) = f(y)$ for all $x, y \in \{0,1\}^n$.
\end{lemma}

\begin{lemma}[Linear Boolean functions have degree at most one]
\lean{LowDegreeTest.linear_is_degree_le_one}
\label{LowDegreeTest.linear_is_degree_le_one}
\leanok
\uses{BoolBLR.is_linear_bool, BoolBLR.lift_pm1, BoolFourier.BoolToPM1, BoolFourier.BoolToPM1_xor, BoolFourier.hypercube, BoolFourier.xor_vec, LowDegreeTest.gowers_product_succ, LowDegreeTest.is_degree_le_bool}
\difficulty{5}
Let $f : \{0,1\}^n \to \{0,1\}$ be linear, meaning $f(x \oplus y) = f(x) \oplus f(y)$
for all $x, y \in \{0,1\}^n$, where $\oplus$ denotes componentwise XOR. Then $f$ has
degree $\le 1$: its $\pm 1$ lift $(-1)^f$ has vanishing multiplicative derivatives of
order two, so that for every base point $x \in \{0,1\}^n$ and every pair of directions
$h_1, h_2 \in \{0,1\}^n$ the order-$2$ Gowers product of $(-1)^f$ over the combinatorial
square $x + \operatorname{span}\{h_1, h_2\}$ equals $1$.
\end{lemma}

\begin{lemma}[Degree-one Boolean functions are affine]
\lean{LowDegreeTest.degree_le_one_implies_affine}
\label{LowDegreeTest.degree_le_one_implies_affine}
\leanok
\statementsource{boolean-ch06-pseudorandomness-f2}{pdf-6c10bfcc8080-p017-b004}
\uses{BoolBLR.is_linear_bool, BoolBLR.lift_pm1, BoolFourier.BoolToPM1, BoolFourier.hypercube, BoolFourier.xor_vec, BoolFourier.zero_vec, LowDegreeTest.gowers_product_succ, LowDegreeTest.is_degree_le_bool}
\difficulty{6}
Let $f : \{0,1\}^n \to \{0,1\}$ be a Boolean function of degree at most $1$; that is,
its $\pm 1$ lift $x \mapsto (-1)^{f(x)}$ has degree at most $1$, meaning every
second-order multiplicative derivative is trivial: for all $x, h_1, h_2 \in \{0,1\}^n$
the order-$2$ Gowers product of $(-1)^f$ at $x$ in directions $h_1, h_2$ equals $1$.
Then $f$ is affine over $\bbf_2$: either $f$ itself is linear, satisfying $f(x \oplus y)
= f(x) \oplus f(y)$ for all $x, y$, or its pointwise negation $x \mapsto \lnot f(x)$ is
linear.
\end{lemma}

\subsection{Reed--Muller codes}

\begin{definition}[Reed--Muller code]
\lean{LowDegreeTest.ReedMuller}
\label{LowDegreeTest.ReedMuller}
\leanok
\uses{BoolFourier.hypercube, LowDegreeTest.is_degree_le_bool}
The Reed--Muller code $\mathrm{RM}(d, m)$ is the set of all Boolean functions
$f : \{0,1\}^m \to \{0,1\}$ of degree $\le d$:
\[
  \mathrm{RM}(d, m) = \{ f : \{0,1\}^m \to \{0,1\} \mid \deg(f) \le d \}.
\]
\end{definition}

\begin{lemma}[Monotonicity of Reed–Muller codes in the degree]
\lean{LowDegreeTest.ReedMuller_monotone}
\label{LowDegreeTest.ReedMuller_monotone}
\leanok
\uses{LowDegreeTest.ReedMuller, LowDegreeTest.degree_le_succ}
\difficulty{1}
Fix an ambient dimension $n$ and a degree $d$. Every Boolean function $f : \{0,1\}^n \to
\{0,1\}$ of degree at most $d$ also has degree at most $d+1$; equivalently, the
Reed–Muller code $\mathrm{RM}(d,n)$ is contained in $\mathrm{RM}(d+1,n)$.
\end{lemma}

\begin{lemma}[The zero function has degree at most $d$]
\lean{LowDegreeTest.zero_mem_ReedMuller}
\label{LowDegreeTest.zero_mem_ReedMuller}
\leanok
\uses{BoolFourier.hypercube, LowDegreeTest.ReedMuller}
\difficulty{3}
Fix $n, d \in \bbn$. The identically-false function $\mathbf{0} : \{0,1\}^n \to
\{0,1\}$, sending every input to $0$, lies in the Reed--Muller code $\mathrm{RM}(d, n)$;
that is, it has degree at most $d$.
\end{lemma}

\begin{lemma}[The constant-true function is low degree]
\lean{LowDegreeTest.one_mem_ReedMuller}
\label{LowDegreeTest.one_mem_ReedMuller}
\leanok
\uses{BoolBLR.lift_pm1, BoolFourier.BoolToPM1, BoolFourier.hypercube, BoolFourier.xor_vec, LowDegreeTest.ReedMuller, LowDegreeTest.gowers_product}
\difficulty{3}
For every $n$ and every degree bound $d$, the constant function $\mathbf{1} : \{0,1\}^n
\to \{0,1\}$ taking the value $1$ at every point lies in the Reed--Muller code
$\mathrm{RM}(d, n)$; that is, it has degree $\le d$.
\end{lemma}

\begin{lemma}[Linear Boolean functions lie in the first-order Reed–Muller code]
\lean{LowDegreeTest.linear_mem_ReedMuller_one}
\label{LowDegreeTest.linear_mem_ReedMuller_one}
\leanok
\statementsource{boolean-ch06-pseudorandomness-f2}{pdf-6c10bfcc8080-p006-b004}
\uses{BoolBLR.is_linear_bool, BoolFourier.hypercube, LowDegreeTest.ReedMuller, LowDegreeTest.linear_is_degree_le_one}
\difficulty{1}
Let $f : \{0,1\}^n \to \{0,1\}$ be a linear Boolean function, meaning that $f(x \oplus
y) = f(x) \oplus f(y)$ for all $x, y \in \{0,1\}^n$, where $\oplus$ denotes
componentwise XOR. Then $f$ belongs to the first-order Reed–Muller code $\mathrm{RM}(1,
n)$; that is, $f$ has degree at most $1$.
\end{lemma}

\begin{lemma}[Degree-one Reed–Muller functions are affine]
\lean{LowDegreeTest.ReedMuller_one_is_affine}
\label{LowDegreeTest.ReedMuller_one_is_affine}
\leanok
\statementsource{boolean-ch06-pseudorandomness-f2}{pdf-6c10bfcc8080-p017-b004}
\uses{BoolBLR.is_linear_bool, BoolFourier.hypercube, LowDegreeTest.ReedMuller, LowDegreeTest.degree_le_one_implies_affine}
\difficulty{1}
Let $f : \{0,1\}^n \to \{0,1\}$ be a Boolean function lying in the Reed–Muller code
$\mathrm{RM}(1,n)$, that is, a function of degree at most $1$. Then $f$ is affine:
either $f$ is $\mathbb{F}_2$-linear, meaning $f(x \oplus y) = f(x) \oplus f(y)$ for all
$x, y \in \{0,1\}^n$, or the complementary function $x \mapsto \lnot f(x)$ is
$\mathbb{F}_2$-linear.
\end{lemma}

\subsection{Gowers norms}

\begin{definition}[Gowers norm (power form)]
\lean{LowDegreeTest.gowers_norm_pow}
\label{LowDegreeTest.gowers_norm_pow}
\leanok
\statementsource{boolean-ch06-pseudorandomness-f2}{pdf-6c10bfcc8080-p026-b004}
\uses{BoolFourier.BoolFun, BoolFourier.expectation, BoolFourier.hypercube, LowDegreeTest.gowers_product, LowDegreeTest.multi_expectation}
The $k$-th Gowers uniformity norm to the $2^k$-th power is
\[
  \|f\|_{U^k}^{2^k}
  = \E_{x \in \{0,1\}^n}\,\E_{h_1,\ldots,h_k}\!\left[
      \mathrm{GP}(f, k, x, h_1,\ldots,h_k)\right].
\]
\end{definition}

\begin{lemma}[Gowers norm one iff bounded degree]
\lean{LowDegreeTest.gowers_norm_eq_one_iff}
\label{LowDegreeTest.gowers_norm_eq_one_iff}
\leanok
\uses{BoolBLR.lift_pm1, BoolFourier.card_hypercube, BoolFourier.expectation, BoolFourier.hypercube, LowDegreeTest.card_multi_hypercube, LowDegreeTest.gowers_norm_pow, LowDegreeTest.gowers_product, LowDegreeTest.gowers_product_pm1, LowDegreeTest.is_degree_le_bool, LowDegreeTest.is_degree_le_pm1, LowDegreeTest.multi_expectation}
\difficulty{6}
Let $f : \{0,1\}^n \to \{0,1\}$ be a Boolean function and let $d$ be a nonnegative
integer, and write $(-1)^f : \{0,1\}^n \to \{\pm 1\}$ for its sign lift. Then the
$2^{d+1}$-th power of the order-$(d+1)$ Gowers uniformity norm of $(-1)^f$ equals $1$,
\[
  \|(-1)^f\|_{U^{d+1}}^{2^{d+1}} = 1,
\]
if and only if $(-1)^f$ has degree at most $d$, meaning that for every base point $x \in
\{0,1\}^n$ and all direction vectors $h_1,\ldots,h_{d+1} \in \{0,1\}^n$ the $(d+1)$-fold
Gowers product $\mathrm{GP}\big((-1)^f, d{+}1, x, h_1,\ldots,h_{d+1}\big)$ equals $1$.
\end{lemma}

\begin{lemma}[The fourth power of the $U^2$ norm as the squared $L^2$ norm of a convolution]
\lean{LowDegreeTest.gowers_U2_eq_L2_conv}
\label{LowDegreeTest.gowers_U2_eq_L2_conv}
\leanok
\uses{BoolFourier.BoolFun, BoolFourier.L2_norm_sq, BoolFourier.convolution, BoolFourier.expectation, BoolFourier.hypercube, BoolFourier.xor_vec, LowDegreeTest.char_is_degree_le_one, LowDegreeTest.gowers_norm_pow, LowDegreeTest.gowers_product, LowDegreeTest.multi_expectation, LowDegreeTest.xor_vec_comm}
\difficulty{6}
Let $f : \{0,1\}^n \to \bbr$ be a Boolean function, and write $x \oplus y$ for the
componentwise XOR of $x, y \in \{0,1\}^n$. Then the fourth power of the second Gowers
uniformity norm of $f$ equals the squared $L^2$ norm of the self-convolution $f * f$,
where $(f * f)(x) = \E_{y}[f(y)\,f(x \oplus y)]$; that is,
\[
  \|f\|_{U^2}^4 = \|f * f\|_2^2 = \E_{x}\bigl[(f * f)(x)^2\bigr].
\]
\end{lemma}

\begin{lemma}[The $U^2$ norm as the $\ell^4$ norm of the Fourier spectrum]
\lean{LowDegreeTest.gowers_U2_fourier}
\label{LowDegreeTest.gowers_U2_fourier}
\leanok
\statementsource{boolean-ch06-pseudorandomness-f2}{pdf-6c10bfcc8080-p003-b001}
\uses{BoolFourier.BoolFun, BoolFourier.fourier_coeff, BoolFourier.fourier_coeff_convolution, BoolFourier.parseval_identity, LowDegreeTest.gowers_U2_eq_L2_conv, LowDegreeTest.gowers_norm_pow}
\difficulty{2}
Let $f : \{0,1\}^n \to \bbr$ be a Boolean function, with Gowers uniformity norm
$\|f\|_{U^2}$ and Fourier coefficients $\hat f(S) = \E_x[f(x)\chi_S(x)]$ indexed by the
subsets $S \subseteq [n]$. Then
\[
  \|f\|_{U^2}^4 \;=\; \sum_{S \subseteq [n]} \hat f(S)^4 .
\]
\end{lemma}

\begin{lemma}[Gowers norm at most one for sign-valued functions]
\lean{LowDegreeTest.gowers_norm_le_one}
\label{LowDegreeTest.gowers_norm_le_one}
\leanok
\uses{BoolBLR.lift_pm1, BoolFourier.card_hypercube, BoolFourier.hypercube, LowDegreeTest.card_multi_hypercube, LowDegreeTest.gowers_norm_pow, LowDegreeTest.gowers_product, LowDegreeTest.gowers_product_pm1}
\difficulty{4}
Let $f : \{0,1\}^n \to \{0,1\}$ be a Boolean-valued function on the hypercube, and let
$g = (-1)^f$ be its $\pm 1$ lift, the real-valued function $g(x) = (-1)^{f(x)} \in \{\pm
1\}$. Then for every $k \in \bbn$, the $k$-th Gowers uniformity norm of $g$ (in power
form) satisfies
\[
  \|g\|_{U^k}^{2^k}
= \E_{x \in \{0,1\}^n}\,\E_{h_1,\ldots,h_k}\!\left[\mathrm{GP}(g, k, x,
h_1,\ldots,h_k)\right]
  \;\le\; 1.
\]
\end{lemma}

\subsection{Gowers product algebra}

\begin{definition}[Pointwise product of Boolean functions]
\lean{LowDegreeTest.boolFun_mul}
\label{LowDegreeTest.boolFun_mul}
\leanok
\uses{BoolFourier.BoolFun}
The pointwise product of two real-valued functions on the hypercube:
$(f \cdot g)(x) = f(x) \cdot g(x)$.
\end{definition}

\begin{lemma}[Multiplicativity of the Gowers product]
\lean{LowDegreeTest.gowers_product_mul}
\label{LowDegreeTest.gowers_product_mul}
\leanok
\uses{BoolFourier.BoolFun, BoolFourier.hypercube, LowDegreeTest.boolFun_mul, LowDegreeTest.gowers_product}
\difficulty{3}
Let $f, g : \{0,1\}^n \to \bbr$ be Boolean functions, fix an order $k \ge 0$, a base
point $x \in \{0,1\}^n$, and direction vectors $h_1, \ldots, h_k \in \{0,1\}^n$, and
write $f \cdot g$ for the pointwise product $(f \cdot g)(y) = f(y)\,g(y)$. Then the
order-$k$ Gowers product of the product function equals the product of the two Gowers
products:
\[
  \mathrm{GP}(f \cdot g,\, k,\, x,\, h_1, \ldots, h_k)
  = \mathrm{GP}(f,\, k,\, x,\, h_1, \ldots, h_k)
    \cdot \mathrm{GP}(g,\, k,\, x,\, h_1, \ldots, h_k).
\]
\end{lemma}

\begin{lemma}[Low-degree factors are absorbed by the Gowers product]
\lean{LowDegreeTest.gowers_product_mul_degree_le}
\label{LowDegreeTest.gowers_product_mul_degree_le}
\leanok
\uses{BoolFourier.BoolFun, BoolFourier.hypercube, LowDegreeTest.boolFun_mul, LowDegreeTest.gowers_product, LowDegreeTest.gowers_product_mul, LowDegreeTest.is_degree_le_pm1}
\difficulty{1}
Let $f, g : \{0,1\}^n \to \bbr$ be Boolean functions, and suppose $g$ has degree $\le
d$, meaning that for every base point $y \in \{0,1\}^n$ and all direction vectors
$h_1,\ldots,h_{d+1} \in \{0,1\}^n$ one has $\mathrm{GP}(g, d{+}1, y, h_1,\ldots,h_{d+1})
= 1$. Then for every base point $x \in \{0,1\}^n$ and all directions $h_1,\ldots,h_{d+1}
\in \{0,1\}^n$, the order-$(d{+}1)$ Gowers product of the pointwise product $f \cdot g$
agrees with that of $f$ alone:
\[
  \mathrm{GP}(f \cdot g,\, d{+}1,\, x,\, h_1,\ldots,h_{d+1})
  = \mathrm{GP}(f,\, d{+}1,\, x,\, h_1,\ldots,h_{d+1}).
\]
\end{lemma}

\begin{lemma}[Gowers norm invariant under multiplication by a low-degree function]
\lean{LowDegreeTest.gowers_norm_mul_degree_le}
\label{LowDegreeTest.gowers_norm_mul_degree_le}
\leanok
\uses{BoolFourier.BoolFun, LowDegreeTest.boolFun_mul, LowDegreeTest.gowers_norm_pow, LowDegreeTest.gowers_product_mul_degree_le, LowDegreeTest.is_degree_le_pm1}
\difficulty{2}
Fix $n, d \in \bbn$ and let $f, g : \{0,1\}^n \to \bbr$ be Boolean functions. If $g$ has
degree $\le d$—that is, every $(d{+}1)$-fold multiplicative derivative of $g$ equals
$1$—then multiplying by $g$ leaves the $(d{+}1)$-st Gowers uniformity norm unchanged:
writing $f \cdot g$ for the pointwise product $(f \cdot g)(x) = f(x)\,g(x)$,
\[
  \|f \cdot g\|_{U^{d+1}}^{2^{d+1}} = \|f\|_{U^{d+1}}^{2^{d+1}}.
\]
\end{lemma}

\subsection{Degree test}

\begin{definition}[Degree test acceptance probability]
\lean{LowDegreeTest.degree_test_accept_prob}
\label{LowDegreeTest.degree_test_accept_prob}
\leanok
\uses{BoolBLR.lift_pm1, BoolFourier.expectation, BoolFourier.hypercube, LowDegreeTest.gowers_product, LowDegreeTest.multi_expectation}
The probability that the $(d{+}1)$-fold derivative test accepts $f$:
\[
  \Pr[\text{accept}] =
  \E_{x}\,\E_{h_1,\ldots,h_{d+1}}
  \bigl[\mathbf{1}[\mathrm{GP}((-1)^f,\, d{+}1,\, x,\, h_1,\ldots,h_{d+1}) = 1]\bigr].
\]
\end{definition}

\begin{definition}[$\varepsilon$-far from degree $d$]
\lean{LowDegreeTest.epsilon_far_from_degree}
\label{LowDegreeTest.epsilon_far_from_degree}
\leanok
\uses{BoolBLR.bool_dist, BoolFourier.hypercube, LowDegreeTest.is_degree_le_bool}
A function $f : \{0,1\}^n \to \{0,1\}$ is $\varepsilon$-far from degree $\le d$
if $0 \le \varepsilon \le 1$ and $\dist(f, g) \ge \varepsilon$ for every Boolean
function $g$ of degree $\le d$.
\end{definition}

\begin{lemma}[Gowers norm and the low-degree test acceptance probability]
\lean{LowDegreeTest.degree_test_accept_prob_eq}
\label{LowDegreeTest.degree_test_accept_prob_eq}
\leanok
\uses{BoolBLR.lift_pm1, BoolFourier.card_hypercube, BoolFourier.expectation, BoolFourier.hypercube, LowDegreeTest.degree_test_accept_prob, LowDegreeTest.gowers_norm_pow, LowDegreeTest.gowers_product, LowDegreeTest.gowers_product_pm1, LowDegreeTest.multi_expectation}
\difficulty{5}
Fix $n$ and an integer $d \ge 0$, and let $f : \{0,1\}^n \to \{0,1\}$ be a
Boolean-valued function with $\pm1$ lift $(-1)^f : \{0,1\}^n \to \{\pm1\}$. Then the
probability that the $(d+1)$-fold derivative low-degree test accepts $f$ is
\[
  \Pr[\text{accept}] \;=\; \frac{1 + \|(-1)^f\|_{U^{d+1}}^{2^{d+1}}}{2},
\]
where $\|(-1)^f\|_{U^{d+1}}^{2^{d+1}}$ is the $(d+1)$-st Gowers uniformity norm of the
lift raised to the power $2^{d+1}$.
\end{lemma}

\begin{lemma}[Completeness of the low-degree test]
\lean{LowDegreeTest.degree_test_completeness}
\label{LowDegreeTest.degree_test_completeness}
\leanok
\uses{BoolFourier.hypercube, LowDegreeTest.degree_test_accept_prob, LowDegreeTest.degree_test_accept_prob_eq, LowDegreeTest.gowers_norm_eq_one_iff, LowDegreeTest.is_degree_le_bool}
\difficulty{2}
Let $f : \{0,1\}^n \to \{0,1\}$ be a Boolean function of degree at most $d$, meaning
that its $\pm 1$ lift $(-1)^f$ has every $(d{+}1)$-fold Gowers product equal to $1$: for
all base points $x$ and all direction vectors $h_1,\ldots,h_{d+1} \in \{0,1\}^n$, one
has $\mathrm{GP}\bigl((-1)^f, d{+}1, x, h_1,\ldots,h_{d+1}\bigr) = 1$. Then the
$(d{+}1)$-fold derivative test accepts $f$ with probability exactly $1$; that is,
\[
\E_{x}\,\E_{h_1,\ldots,h_{d+1}}\bigl[\mathbf{1}[\mathrm{GP}((-1)^f,\, d{+}1,\, x,\,
h_1,\ldots,h_{d+1}) = 1]\bigr] = 1.
\]
\end{lemma}

\begin{lemma}[Qualitative soundness of the degree test]
\lean{LowDegreeTest.degree_test_qualitative_soundness}
\label{LowDegreeTest.degree_test_qualitative_soundness}
\leanok
\uses{BoolFourier.hypercube, BoolFourier.xor_vec, LowDegreeTest.degree_test_accept_prob, LowDegreeTest.degree_test_accept_prob_eq, LowDegreeTest.gowers_norm_eq_one_iff, LowDegreeTest.gowers_norm_pow, LowDegreeTest.gowers_product_pm1, LowDegreeTest.is_degree_le_bool}
\difficulty{5}
Let $f : \{0,1\}^n \to \{0,1\}$ be a Boolean function that does not have degree at most
$d$, meaning that its $\pm 1$ lift $(-1)^f$ has at least one nonvanishing $(d{+}1)$-fold
multiplicative derivative: there exist a base point $x$ and directions
$h_1,\ldots,h_{d+1} \in \{0,1\}^n$ with $\mathrm{GP}\!\left((-1)^f,\, d{+}1,\, x,\,
h_1,\ldots,h_{d+1}\right) \neq 1$. Then the $(d{+}1)$-fold derivative test accepts $f$
with probability strictly less than $1$.
\end{lemma}

\subsection{Fourier analysis and degree}

\begin{lemma}[Walsh characters have degree at most one]
\lean{LowDegreeTest.char_is_degree_le_one}
\label{LowDegreeTest.char_is_degree_le_one}
\leanok
\statementsource{boolean-ch06-pseudorandomness-f2}{pdf-6c10bfcc8080-p006-b004}
\uses{BoolFourier.char_S, BoolFourier.char_S_sq, BoolFourier.hypercube, BoolFourier.xor_vec, BooleanAnalysis.boolToSign, BooleanAnalysis.chiS, LowDegreeTest.gowers_product, LowDegreeTest.is_degree_le_pm1}
\difficulty{4}
For every subset $S \subseteq [n]$, the Walsh--Fourier character $\chi_S$ has degree at
most $1$: every second-order Gowers product of $\chi_S$ equals $1$. Explicitly, for all
base points $x \in \{0,1\}^n$ and all direction vectors $h_1, h_2 \in \{0,1\}^n$,
\[
\chi_S(x)\,\chi_S(x \oplus h_1)\,\chi_S(x \oplus h_2)\,\chi_S(x \oplus h_1 \oplus h_2)
\;=\; 1.
\]
\end{lemma}

\begin{lemma}[Negated Walsh character has degree at most one]
\lean{LowDegreeTest.neg_char_is_degree_le_one}
\label{LowDegreeTest.neg_char_is_degree_le_one}
\leanok
\statementsource{boolean-ch06-pseudorandomness-f2}{pdf-6c10bfcc8080-p006-b004}
\uses{BoolFourier.char_S, LowDegreeTest.char_is_degree_le_one, LowDegreeTest.is_degree_le_pm1}
\difficulty{2}
For every subset $S \subseteq [n]$, the negated Walsh--Fourier character $x \mapsto
-\chi_S(x)$ on $\{0,1\}^n$ has degree $\le 1$: its second-order multiplicative
derivatives all equal $1$. Explicitly, for every base point $x \in \{0,1\}^n$ and every
pair of direction vectors $h_1, h_2 \in \{0,1\}^n$, the order-$2$ Gowers product of
$-\chi_S$ satisfies $\mathrm{GP}(-\chi_S, 2, x, h_1, h_2) = 1$.
\end{lemma}

\begin{lemma}[Walsh characters have bounded multiplicative degree]
\lean{LowDegreeTest.char_is_degree_le}
\label{LowDegreeTest.char_is_degree_le}
\leanok
\uses{BoolFourier.char_S, LowDegreeTest.char_is_degree_le_one, LowDegreeTest.degree_le_succ, LowDegreeTest.is_degree_le_pm1}
\difficulty{3}
For every subset $S \subseteq [n]$ and every integer $d \ge 1$, the Walsh--Fourier
character $\chi_S$ has degree $\le d$: for every base point $x \in \{0,1\}^n$ and all
direction vectors $h_1, \ldots, h_{d+1} \in \{0,1\}^n$, the $(d{+}1)$-fold Gowers
product satisfies $\mathrm{GP}(\chi_S,\, d{+}1,\, x,\, h_1, \ldots, h_{d+1}) = 1$.
\end{lemma}

\begin{lemma}[Negated Walsh character has degree at most $d$]
\lean{LowDegreeTest.neg_char_is_degree_le}
\label{LowDegreeTest.neg_char_is_degree_le}
\leanok
\uses{BoolFourier.char_S, LowDegreeTest.degree_le_succ, LowDegreeTest.is_degree_le_pm1, LowDegreeTest.neg_char_is_degree_le_one}
\difficulty{3}
For every subset $S \subseteq [n]$ and every integer $d \ge 1$, the negated Walsh
character $x \mapsto -\chi_S(x)$ has degree $\le d$. Equivalently, all of its
$(d{+}1)$-fold multiplicative derivatives are trivial: for every base point $x \in
\{0,1\}^n$ and all direction vectors $h_1,\ldots,h_{d+1} \in \{0,1\}^n$, the
order-$(d{+}1)$ Gowers product satisfies $\mathrm{GP}(-\chi_S,\, d{+}1,\, x,\,
h_1,\ldots,h_{d+1}) = 1$.
\end{lemma}

\begin{lemma}[Fourier coefficients of a function far from low degree]
\lean{LowDegreeTest.fourier_coeff_le_of_far_from_degree}
\label{LowDegreeTest.fourier_coeff_le_of_far_from_degree}
\leanok
\uses{BoolBLR.fourier_coeff_le_of_dist_ge, BoolBLR.lift_pm1, BoolFourier.char_S, BoolFourier.fourier_coeff, BoolFourier.hypercube, BooleanAnalysis.chiS, LowDegreeTest.char_is_degree_le, LowDegreeTest.epsilon_far_from_degree, LowDegreeTest.is_degree_le_bool}
\difficulty{5}
Let $f:\{0,1\}^n\to\{0,1\}$ be a Boolean function, let $d\ge 1$, and let $\varepsilon$
satisfy $0\le\varepsilon\le 1$. Suppose $f$ is $\varepsilon$-far from degree $\le d$;
that is, $\Pr_{x}[f(x)\ne g(x)]\ge\varepsilon$ for every Boolean function
$g:\{0,1\}^n\to\{0,1\}$ of degree $\le d$. Writing $(-1)^f:\{0,1\}^n\to\{\pm1\}$ for the
$\pm1$ lift of $f$ and $\widehat{(-1)^f}(S)=\E_x[(-1)^{f(x)}\chi_S(x)]$ for its Fourier
coefficient at $S$, we then have, for every $S\subseteq[n]$,
\[
  \widehat{(-1)^f}(S)\;\le\;1-2\varepsilon.
\]
\end{lemma}

\begin{lemma}[Fourier coefficients bounded below when far from low degree]
\lean{LowDegreeTest.neg_fourier_coeff_le_of_far_from_degree}
\label{LowDegreeTest.neg_fourier_coeff_le_of_far_from_degree}
\leanok
\uses{BoolBLR.bool_dist, BoolBLR.lift_pm1, BoolFourier.BoolToPM1, BoolFourier.PM1ToBool?, BoolFourier.char_S, BoolFourier.expectation, BoolFourier.fourier_coeff, BoolFourier.hypercube, BoolFourier.inner_product, LowDegreeTest.epsilon_far_from_degree, LowDegreeTest.is_degree_le_bool, LowDegreeTest.is_degree_le_pm1, LowDegreeTest.neg_char_is_degree_le}
\difficulty{6}
Let $f : \{0,1\}^n \to \{0,1\}$ be a Boolean function, let $d \ge 1$, and write $F =
(-1)^f : \{0,1\}^n \to \{\pm 1\}$ for its $\pm 1$ lift. Suppose $f$ is $\varepsilon$-far
from degree $\le d$; that is, $0 \le \varepsilon \le 1$ and every Boolean function $g$
of degree $\le d$ satisfies $\Pr_{x}[f(x) \ne g(x)] \ge \varepsilon$. Then for every $S
\subseteq [n]$,
\[
  -\widehat{F}(S) \;\le\; 1 - 2\varepsilon,
\]
where $\widehat{F}(S) = \mathbb{E}_x\bigl[F(x)\,\chi_S(x)\bigr]$ is the Fourier
coefficient of $F$ at $S$; equivalently, $\widehat{F}(S) \ge 2\varepsilon - 1$.
\end{lemma}

\begin{lemma}[Fourier coefficients of a function far from low degree]
\lean{LowDegreeTest.abs_fourier_coeff_le_of_far_from_degree}
\label{LowDegreeTest.abs_fourier_coeff_le_of_far_from_degree}
\leanok
\uses{BoolBLR.lift_pm1, BoolFourier.fourier_coeff, BoolFourier.hypercube, LowDegreeTest.epsilon_far_from_RM, LowDegreeTest.epsilon_far_from_degree, LowDegreeTest.fourier_coeff_le_of_far_from_degree, LowDegreeTest.neg_fourier_coeff_le_of_far_from_degree}
\difficulty{2}
Let $f : \{0,1\}^n \to \{0,1\}$ be a Boolean function, let $d \ge 1$, and let
$\varepsilon$ be a real number with $0 \le \varepsilon \le 1$. Suppose $f$ is
$\varepsilon$-far from degree $\le d$, meaning that $\dist(f,g) \ge \varepsilon$ for
every Boolean function $g$ of degree $\le d$. Then, writing $\widehat{(-1)^f}(S)$ for
the Fourier coefficient at $S$ of the $\pm 1$ lift $x \mapsto (-1)^{f(x)}$, we have
\[
  \bigl|\widehat{(-1)^f}(S)\bigr| \;\le\; 1 - 2\varepsilon
  \qquad\text{for every } S \subseteq [n].
\]
\end{lemma}

\subsection{Quantitative soundness}

\begin{definition}[$\varepsilon$-far from Reed--Muller]
\lean{LowDegreeTest.epsilon_far_from_RM}
\label{LowDegreeTest.epsilon_far_from_RM}
\leanok
\uses{BoolFourier.hypercube, LowDegreeTest.epsilon_far_from_degree}
A function $f$ is $\varepsilon$-far from $\mathrm{RM}(d, n)$ if it is $\varepsilon$-far
from degree $\le d$ in the sense of
\texttt{LowDegreeTest.epsilon\_far\_from\_degree}.
\end{definition}

\begin{definition}[Reed--Muller test acceptance probability]
\lean{LowDegreeTest.RM_test_accept_prob}
\label{LowDegreeTest.RM_test_accept_prob}
\leanok
\uses{BoolFourier.hypercube, LowDegreeTest.degree_test_accept_prob}
The acceptance probability of the Reed--Muller test, defined as the degree-test
acceptance probability \texttt{LowDegreeTest.degree\_test\_accept\_prob}.
\end{definition}

\begin{lemma}[Parseval identity for the sign lift of a Boolean function]
\lean{LowDegreeTest.parseval_pm1}
\label{LowDegreeTest.parseval_pm1}
\leanok
\statementsource{boolean-ch03-spectral-learning-restrictions}{
  pdf-20e15a3d3d5e-p004-b001,
  pdf-20e15a3d3d5e-p001-b003
}
\uses{BoolBLR.lift_pm1, BoolFourier.BoolToPM1, BoolFourier.BoolToPM1_sq, BoolFourier.L2_norm_sq, BoolFourier.expectation, BoolFourier.fourier_coeff, BoolFourier.hypercube, BoolFourier.parseval_identity}
\difficulty{2}
Let $f : \{0,1\}^n \to \{0,1\}$ be a Boolean-valued function, and let $g : \{0,1\}^n \to
\bbr$ be its $\pm 1$ lift $g(x) = (-1)^{f(x)}$. Writing $\hat g(S) = \E[g\,\chi_S]$ for
the Fourier coefficient of $g$ at $S$, where $\chi_S(x) = \prod_{i \in S}(-1)^{x_i}$ and
the expectation is over $x$ uniform on $\{0,1\}^n$, one has
\[
  \sum_{S \subseteq [n]} \hat g(S)^2 = 1.
\]
\end{lemma}

\begin{lemma}[$U^2$ norm bound for functions far from low degree]
\lean{LowDegreeTest.gowers_U2_le_of_far}
\label{LowDegreeTest.gowers_U2_le_of_far}
\leanok
\uses{BoolBLR.lift_pm1, BoolFourier.fourier_coeff, BoolFourier.hypercube, LowDegreeTest.abs_fourier_coeff_le_of_far_from_degree, LowDegreeTest.epsilon_far_from_degree, LowDegreeTest.gowers_U2_fourier, LowDegreeTest.gowers_norm_pow, LowDegreeTest.parseval_pm1}
\difficulty{4}
Let $f : \{0,1\}^n \to \{0,1\}$ be a Boolean function, let $d \ge 1$, and let
$\varepsilon$ satisfy $0 \le \varepsilon \le 1$. Suppose $f$ is $\varepsilon$-far from
degree $\le d$; that is, $\dist(f,g) \ge \varepsilon$ for every Boolean function $g$ of
degree $\le d$. Then the $\pm 1$ lift $(-1)^f : \{0,1\}^n \to \{\pm 1\}$ satisfies \[
\bigl\lVert (-1)^f \bigr\rVert_{U^2}^{4} \;\le\; (1 - 2\varepsilon)^2, \] where
$\lVert\,\cdot\,\rVert_{U^2}^{4}$ denotes the second Gowers uniformity norm raised to
the fourth power,
$\E_{x}\,\E_{h_1,h_2}\!\bigl[\mathrm{GP}\bigl((-1)^f,2,x,h_1,h_2\bigr)\bigr]$.
\end{lemma}

\begin{lemma}[Squaring contracts on the unit interval]
\lean{LowDegreeTest.sq_one_sub_two_eps_le}
\label{LowDegreeTest.sq_one_sub_two_eps_le}
\leanok
\difficulty{2}
For every real number $\varepsilon$ with $0 \le \varepsilon \le \tfrac{1}{2}$, the
quantity $1 - 2\varepsilon$ lies in $[0,1]$ and hence is not increased by squaring:
\[
(1 - 2\varepsilon)^2 \le 1 - 2\varepsilon.
\]
\end{lemma}

\begin{lemma}[Farness from low degree is at most one half]
\lean{LowDegreeTest.eps_le_half_of_far}
\label{LowDegreeTest.eps_le_half_of_far}
\leanok
\uses{BoolBLR.lift_pm1, BoolFourier.fourier_coeff, BoolFourier.hypercube, LowDegreeTest.abs_fourier_coeff_le_of_far_from_degree, LowDegreeTest.epsilon_far_from_degree}
\difficulty{3}
Fix $n$ and an integer $d \ge 1$, and let $f : \{0,1\}^n \to \{0,1\}$ be a Boolean
function. If $f$ is $\varepsilon$-far from degree $\le d$ — that is, $0 \le \varepsilon
\le 1$ and every Boolean function $g$ of degree $\le d$ satisfies $\dist(f,g) \ge
\varepsilon$ — then $\varepsilon \le 1/2$.
\end{lemma}

\begin{lemma}[Gowers $U^2$ norm bound for functions far from low degree]
\lean{LowDegreeTest.gowers_norm_le_of_far_d1}
\label{LowDegreeTest.gowers_norm_le_of_far_d1}
\leanok
\uses{BoolBLR.lift_pm1, BoolFourier.hypercube, LowDegreeTest.eps_le_half_of_far, LowDegreeTest.epsilon_far_from_degree, LowDegreeTest.gowers_U2_le_of_far, LowDegreeTest.gowers_norm_pow, LowDegreeTest.sq_one_sub_two_eps_le}
\difficulty{2}
Let $d \ge 1$ and let $f : \{0,1\}^n \to \{0,1\}$ be a Boolean function, and write $F =
(-1)^f$ for its $\pm 1$ lift. Suppose $f$ is $\varepsilon$-far from degree $\le d$: that
is, $0 \le \varepsilon \le 1$, and every Boolean function $g$ of degree $\le d$
satisfies $\dist(f,g) \ge \varepsilon$, where $\dist(f,g) = \Pr_{x}[f(x) \ne g(x)]$ over
a uniform $x \in \{0,1\}^n$. Then the second Gowers uniformity norm of $F$ obeys
\[
  \|F\|_{U^2}^{4} \;\le\; 1 - 2\varepsilon .
\]
\end{lemma}

\begin{lemma}[Monotonicity of $\varepsilon$-farness in the degree bound]
\lean{LowDegreeTest.epsilon_far_monotone}
\label{LowDegreeTest.epsilon_far_monotone}
\leanok
\uses{BoolBLR.lift_pm1, BoolFourier.hypercube, LowDegreeTest.degree_le_succ, LowDegreeTest.epsilon_far_from_degree, LowDegreeTest.is_degree_le_pm1}
\difficulty{3}
Fix a Boolean function $f : \{0,1\}^n \to \{0,1\}$, a real number $\varepsilon$ with $0
\le \varepsilon \le 1$, and natural numbers $d' \le d$. If $f$ is $\varepsilon$-far from
degree $\le d$ — that is, $\dist(f, g) \ge \varepsilon$ for every Boolean function $g$
of degree $\le d$ — then $f$ is also $\varepsilon$-far from degree $\le d'$.
\end{lemma}

\begin{lemma}[Distance of the discrete derivative from low degree]
\lean{LowDegreeTest.derivative_distance_lemma}
\label{LowDegreeTest.derivative_distance_lemma}
\leanok
\uses{BoolBLR.bool_dist, BoolFourier.expectation, BoolFourier.hypercube, BoolFourier.xor_vec, LowDegreeTest.epsilon_far_from_degree, LowDegreeTest.is_degree_le_bool}
\difficulty{4}
Let $f : \{0,1\}^n \to \{0,1\}$ be a Boolean function that is $\varepsilon$-far from
degree $\le d$, where $d \ge 2$; in particular $0 \le \varepsilon \le 1$. For a shift $h
\in \{0,1\}^n$, write $\partial_h f$ for the discrete derivative $\partial_h f(x) = f(x)
\oplus f(x \oplus h)$. Then there exists a real number $\delta \ge \varepsilon$ such
that for every $h \in \{0,1\}^n$ there is a Boolean function $g : \{0,1\}^n \to \{0,1\}$
for which, if $g$ has degree $\le d-1$, then $\dist(\partial_h f, g) \ge \delta$.
\end{lemma}

\begin{lemma}[Gowers norm bound for functions far from low degree]
\lean{LowDegreeTest.gowers_norm_le_of_far}
\label{LowDegreeTest.gowers_norm_le_of_far}
\leanok
\uses{BoolBLR.lift_pm1, BoolFourier.hypercube, LowDegreeTest.epsilon_far_from_degree, LowDegreeTest.gowers_norm_le_of_far_d1, LowDegreeTest.gowers_norm_pow}
\difficulty{3}
Let $d \ge 1$ and let $f : \{0,1\}^n \to \{0,1\}$ be a Boolean function that is
$\varepsilon$-far from degree $\le d$, meaning $0 \le \varepsilon \le 1$ and $\dist(f,
g) \ge \varepsilon$ for every Boolean function $g$ of degree $\le d$. Then there is a
constant $c > 0$ for which the $\pm 1$ lift $(-1)^{f}$ satisfies
\[
  \bigl\|(-1)^{f}\bigr\|_{U^{d+1}}^{2^{d+1}} \;\le\; 1 - c\,\varepsilon.
\]
\end{lemma}

\begin{lemma}[Quantitative soundness of the degree test]
\lean{LowDegreeTest.degree_test_quantitative_soundness}
\label{LowDegreeTest.degree_test_quantitative_soundness}
\leanok
\uses{BoolFourier.hypercube, LowDegreeTest.degree_test_accept_prob, LowDegreeTest.epsilon_far_from_degree}
Fix $d \ge 1$ and let $f : \{0,1\}^n \to \{0,1\}$ be a Boolean function that is
$\varepsilon$-far from degree $\le d$; that is, $0 \le \varepsilon \le 1$ and every
Boolean function $g$ of degree $\le d$ satisfies $\dist(f,g) \ge \varepsilon$. Then the
$(d{+}1)$-fold derivative test accepts $f$ with probability at most $1 - \varepsilon$:
\[
  \E_{x}\,\E_{h_1,\ldots,h_{d+1}}
  \bigl[\mathbf{1}[\mathrm{GP}((-1)^f, d{+}1, x, h_1,\ldots,h_{d+1}) = 1]\bigr]
  \;\le\; 1 - \varepsilon,
\]
where $x$ and the direction vectors $h_1,\ldots,h_{d+1}$ range uniformly over
$\{0,1\}^n$.
\end{lemma}

\begin{lemma}[Completeness of the Reed–Muller degree test]
\lean{LowDegreeTest.RM_test_completeness}
\label{LowDegreeTest.RM_test_completeness}
\leanok
\uses{BoolFourier.hypercube, LowDegreeTest.ReedMuller, LowDegreeTest.degree_test_accept_prob, LowDegreeTest.degree_test_completeness}
\difficulty{1}
Fix a degree bound $d$ and let $f : \{0,1\}^n \to \{0,1\}$ be a Boolean function
belonging to the Reed–Muller code $\mathrm{RM}(d,n)$, that is, a function of degree at
most $d$. Then the $(d+1)$-fold derivative test accepts $f$ with probability exactly
$1$:
\[
\E_{x}\,\E_{h_1,\ldots,h_{d+1}}\bigl[\mathbf{1}[\mathrm{GP}((-1)^f,\, d{+}1,\, x,\,
h_1,\ldots,h_{d+1}) = 1]\bigr] = 1,
\]
where $x$ and the direction vectors $h_1,\ldots,h_{d+1}$ range uniformly over
$\{0,1\}^n$.
\end{lemma}

\begin{lemma}[Soundness of the low-degree derivative test]
\lean{LowDegreeTest.RM_test_soundness}
\label{LowDegreeTest.RM_test_soundness}
\leanok
\uses{BoolFourier.hypercube, LowDegreeTest.degree_test_accept_prob, LowDegreeTest.degree_test_quantitative_soundness, LowDegreeTest.epsilon_far_from_degree}
\difficulty{1}
Fix $d \ge 1$ and let $f : \{0,1\}^n \to \{0,1\}$ be a Boolean function. Suppose $f$ is
$\varepsilon$-far from degree $\le d$; that is, $0 \le \varepsilon \le 1$ and
$\dist(f,g) \ge \varepsilon$ for every Boolean function $g$ of degree $\le d$. Then the
$(d{+}1)$-fold derivative test accepts $f$ with probability at most $1 - \varepsilon$.
\end{lemma}
